\begin{scanpage}{0045}{28}
and
\[\left.\begin{aligned}f^0&=\frac{\sin\phi_0}{\cos\phi+\cos\phi_0},\\g^0&=-\frac{\sin\phi}{\cos\phi+\cos\phi_0},\end{aligned}\right\}\tag{4.09}\]
if one face of the wedge is illuminated (that is, when $0<\phi_0<\alpha-\pi$), and
\[\left.\begin{aligned}f^0&=\frac{\sin\phi_0}{\cos\phi+\cos\phi_0}+\frac{\sin(\alpha-\phi_0)}{\cos(\alpha-\phi)+\cos(\alpha-\phi_0)},\\g^0&=-\frac{\sin\phi}{\cos\phi+\cos\phi_0}-\frac{\sin(\alpha-\phi)}{\cos(\alpha-\phi)+\cos(\alpha-\phi_0)},\end{aligned}\right\}\tag{4.10}\]
if both faces of the wedge are illuminated (that is, when $\alpha-\pi<\phi_0<\pi$). Let us recall that the functions $f$ and $g$ describe the cylindrical waves radiated by the total current --- that is, the sum of the uniform and nonuniform parts, and the functions $f^0$ and $g^0$ refer to the cylindrical waves radiated only by the uniform part of the current $(\mathbf j^0)$.

Let us note certain properties of the functions $f^1$ and $g^1$. The function $f^1=f^1(\alpha,\phi,\phi_0)$ is continuous, whereas the function $g^1=g^1(\alpha,\phi,\phi_0)$ undergoes a finite discontinuity when $\phi_0=\alpha-\pi$. The reason for this discontinuity is that the uniform part of the current differs from zero on the face along which plane wave (1.24) is propagated (with $\phi_0=\alpha-\pi$). In the case of radar, when the direction to the observation point coincides with the direction to the source $(\phi=\phi_0)$, both functions $f^1$ and $g^1$ are continuous. There is no discontinuity of the function $g^1$ with $\phi=\phi_0=\alpha-\pi$, because the current element does not radiate in the longitudinal direction.

On the boundary of the plane waves (that is, when $\phi=\pi\pm\phi_0$ and $\phi=2\alpha-\pi-\phi_0$) the functions $f$, $f^0$ and $g$, $g^0$ become infinite, whereas the functions $f^1$ and $g^1$ remain finite. In accordance with Equations (4.07)--(4.10), they take the following values
\end{scanpage}
